\documentclass[11pt]{article}
\usepackage{geometry}
\usepackage{amsmath}
\usepackage{amssymb}
\usepackage{booktabs}
\usepackage{siunitx}
\usepackage{hyperref}
\usepackage{graphicx}
\geometry{margin=1in}

\title{Line-Heating Plate Bending: Governing Equations, Methodology, Test Cases, and Calibrated Parameters}
\author{Project: Ship\_Plate\_Bending\_LineHeating}
\date{20 Jan 2026}

\begin{document}
\maketitle

\begin{abstract}
This report documents the calibrated inherent-strain model used to match the four Li2023 benchmark cases
(case IDs 1, 6, 87, 90) within the line-heating plate-bending workflow. It summarizes the governing
thermo-mechanical equations, numerical methodology, test-case definitions, parametric sweeps, and the
final optimized parameters used in the calibrated runs.
\end{abstract}

\section{Governing Equations}
\subsection{Numerical Solution}
The heat equation is discretized in space using linear tetrahedral finite elements and integrated
in time with an implicit backward-Euler scheme. At each time step, the Gaussian heat source is
updated based on the scan position, and the resulting linear system is solved for the temperature
field. Convection and radiation are applied as boundary flux terms. The mechanical problem is
solved as a linear thermo-elastic system with an inherent strain field applied in the heating
band. The thermal history provides temperature-dependent thermal strain, and the displacement
field is obtained by solving the linear elasticity system for each case after the thermal solve.

\subsection{Heat Transfer}
The transient heat equation is solved on a 3D plate domain:
\begin{equation}
\rho c_p \frac{\partial T}{\partial t} - \nabla \cdot (k \nabla T) = q(\mathbf{x}, t),
\end{equation}
where $\rho$ is density, $c_p$ is specific heat, $k$ is thermal conductivity, and $q$ is the applied heat
source. Temperature-dependent material properties are approximated with linear slopes derived from tabulated
values.

\subsection{Heat Source}
A moving Gaussian surface heat flux is applied along the heating line:
\begin{equation}
q(r) = q_0 \exp\left(-2\frac{r^2}{r_0^2}\right),
\end{equation}
where $r$ is radial distance from the heat centerline and $r_0$ is the Gaussian radius. The peak flux $q_0$
corresponds to the specified energy per unit length $E_{in}$ and scan speed $v$:
\begin{equation}
q_0 = \frac{2 E_{in} v}{\pi r_0^2}.
\end{equation}

\subsection{Boundary Conditions}
Convection and radiation are applied on all surfaces:
\begin{equation}
- k \nabla T \cdot \mathbf{n} = h_{conv} (T - T_\infty) + \varepsilon \sigma (T^4 - T_\infty^4).
\end{equation}
Mechanical boundary conditions use a centerline-fixed support model for the Li2023 cases.

\subsection{Thermo-Mechanical Strain}
The total strain is decomposed into elastic, thermal, and inherent components:
\begin{equation}
\varepsilon = \varepsilon^e + \varepsilon^{th} + \varepsilon^{inh},
\end{equation}
with thermal strain $\varepsilon^{th} = \alpha (T - T_0)$ and inherent strain
$\varepsilon^{inh} = \varepsilon_0 \, \psi(\mathbf{x})$ applied in a band around the heating line.

\section{Problem Description}
Line heating induces localized thermal expansion, followed by cooling and contraction that leaves
residual strains and permanent curvature. The objective is to reproduce experimentally observed
plate deflections for benchmark cases by calibrating an inherent-strain law that maps process
inputs (energy, passes, thickness) to an effective shrinkage strain applied in a narrow band
around the heating line. The workflow balances accuracy in single-pass and multi-pass cases over
both low- and high-energy regimes.

\section{Methodology}
\subsection{Workflow}
\begin{enumerate}
  \item Build a 3D plate mesh with a refined band around the heating line.
  \item Solve transient heat transfer with a moving Gaussian heat flux.
  \item Convert thermal history to inherent strain using a calibrated model.
  \item Solve the mechanical deformation using linear thermo-elasticity with inherent strain.
  \item Report global deflection and compare to experimental values.
\end{enumerate}

\subsection{Numerical Settings (Li2023 Runs)}
\begin{itemize}
  \item Geometry: $L_x=400$ mm, $L_y=300$ mm; thickness varies per case.
  \item Heat radius: $r_0=7$ mm; Gaussian exponent $\beta=3$.
  \item Mesh: $h=20$ mm, $h_{refine}=5$ mm, refine band $=60$ mm.
  \item Time stepping: $\Delta t=2$ s, extra cooling $=200$ s.
  \item Boundary: centerline-fixed.
  \item Inherent band: $\sigma_{inh}=7$ mm, $z_{frac}=0.2$.
\end{itemize}

\subsection{Mesh Description}
The mesh is a 3D tetrahedral discretization of the plate, with a locally refined band aligned
to the heating line. The refined band resolves steep thermal gradients around the moving heat
source and improves accuracy in the resulting curvature field. A coarser mesh is used away from
the band to reduce computational cost while preserving global stiffness and boundary effects.

\subsection{Representative Fields}
When available, the following figures summarize representative fields from the calibrated runs.
\begin{figure}[h]
\centering
\IfFileExists{../results/li2023_case_001_v8_full_four_cases_lowE_single/mesh_3d.png}{
  \includegraphics[width=0.48\textwidth]{../results/li2023_case_001_v8_full_four_cases_lowE_single/mesh_3d.png}
  \includegraphics[width=0.48\textwidth]{../results/li2023_case_001_v8_full_four_cases_lowE_single/temperature_top.png}
}{\fbox{Representative mesh/temperature figures not found.}}
\caption{Representative mesh (left) and top-surface temperature field (right) for case 1.}
\end{figure}

\begin{figure}[h]
\centering
\IfFileExists{../results/li2023_case_001_v8_full_four_cases_lowE_single/deflection_top.png}{
  \includegraphics[width=0.48\textwidth]{../results/li2023_case_001_v8_full_four_cases_lowE_single/deflection_top.png}
  \includegraphics[width=0.48\textwidth]{../results/li2023_case_001_v8_full_four_cases_lowE_single/deflection_3d.png}
}{\fbox{Representative deflection figures not found.}}
\caption{Representative deflection contours for case 1 (top surface and 3D view).}
\end{figure}

\section{Test Cases}
Table~\ref{tab:cases} lists the four benchmark cases.
\begin{table}[h]
\centering
\begin{tabular}{cccccc}
\toprule
Case & $v$ (mm/s) & $E_{in}$ (J/mm) & $H$ (mm) & Passes & $w_{exp}$ (mm) \\
\midrule
1  & 6  & 600  & 14 & 1 & 5.785 \\
6  & 11 & 600  & 14 & 2 & 7.712 \\
87 & 12 & 1000 & 10 & 2 & 13.506 \\
90 & 15 & 1000 & 10 & 1 & 11.322 \\
\bottomrule
\end{tabular}
\caption{Li2023 benchmark cases used for calibration.}
\label{tab:cases}
\end{table}

\section{Parametric Sweeps}
The following parameters were swept during calibration:
\begin{itemize}
  \item Base law exponents: $a$ (energy), $p$ (thickness), and $C$ (scale).
  \item Pass scaling: $\delta$ (per-pass increment), $k_{energy}$ (high-energy attenuation),
        $k_{low}$ (low-energy multi-pass attenuation), and $E_{ref,pass}$.
  \item Single-pass scaling: $k_{single}$ (high-energy single-pass boost).
  \item Low-energy single-pass boost: $\eta$, $E_0$, and $\Delta_E$.
\end{itemize}
Sweeps were evaluated against the four cases, aiming to minimize deflection error while
preserving good fits in high-energy and multi-pass regimes.

\section{Optimized Parameters}
Final calibrated parameters (used in the ``final\_calibrated'' tag):
\begin{align*}
C &= 0.06, & a &= 0.35, & b &= 0.0, & p &= 0.325, \\
\delta &= 0.02, & k_{energy} &= 0.43, & k_{low} &= 0.88, & E_{ref,pass} &= 800, \\
k_{single} &= 0.33, & \eta &= 0.2, & E_0 &= 600, & \Delta_E &= 180.
\end{align*}

\section{Inherent-Strain Model}
The inherent strain magnitude $\varepsilon_0$ is computed as:
\begin{equation}
\varepsilon_0 = C \left(\frac{E}{E_{ref}}\right)^a \left(\frac{H_{ref}}{H}\right)^p
\cdot f_{pass} \cdot f_{single} \cdot f_{low,1},
\end{equation}
where $f_{pass}$ includes multi-pass scaling with energy attenuation, $f_{single}$ boosts
high-energy single-pass cases, and $f_{low,1}$ boosts low-energy single-pass cases:
\begin{align}
f_{pass} &= \left(1 + (N-1)\delta\right)\,\max\left(0, 1 - k_{energy}(E/E_{ref}-1)\right)\,\max\left(0, 1 - k_{low}(1 - E/E_{ref,pass})\right),\\
f_{single} &= 1 + k_{single}\max\left(0, E/E_{ref} - 1\right),\\
f_{low,1} &= 1 + \eta\exp\left(-\left(\frac{E - E_0}{\Delta_E}\right)^2\right)\quad (N=1).
\end{align}

\section{Results}
Table~\ref{tab:results} summarizes the final calibrated predictions.
\begin{table}[h]
\centering
\begin{tabular}{ccccc}
\toprule
Case & $w_{exp}$ (mm) & $w_{pred}$ (mm) & Error (mm) & $\varepsilon_0$ \\
\midrule
1  & 5.785 & 5.888 & +0.103 & 0.0720 \\
6  & 7.712 & 7.722 & +0.010 & 0.047736 \\
87 & 13.506 & 13.486 & -0.020 & 0.058235 \\
90 & 11.322 & 11.319 & -0.003 & 0.097785 \\
\bottomrule
\end{tabular}
\caption{Final calibrated results using the low-energy single-pass boost.}
\label{tab:results}
\end{table}

\subsection{Discussion}
The low-energy single-pass boost corrects underprediction in case 1 without affecting multi-pass
cases (case 6) or high-energy cases (87, 90). The low-energy multi-pass attenuation ensures that
case 6 remains close to the experimental target while the high-energy attenuation preserves the
two-pass behavior for case 87. The resulting fit achieves sub-millimeter error across all four
benchmark cases.

\section{Reproducibility}
All run configurations and outputs are stored under the results folder. The latest
calibrated comparison table is available in
\texttt{results/li2023\_comparison\_v8\_full\_four\_cases\_lowE\_single.csv}.

\section{Planning and Post-Processing Workflow}
To transition from calibrated physics to production use, the workflow can be extended with
an inverse planner and NC post-processing pipeline. The goal is to map a target curvature
profile to a discrete, machine-executable heating plan.

\subsection{Inverse Planner}
Define a low-dimensional parameter vector $p$ that encodes heating line positions, pass counts,
energy per unit length, and scan speeds. Given a desired camber profile $w_{target}(y)$, the
planner minimizes a misfit:
\begin{equation}
J(p) = \int_0^{L_y} \left(w_{sim}(y; p) - w_{target}(y)\right)^2\,dy + \lambda\,\mathcal{P}(p),
\end{equation}
where $\mathcal{P}(p)$ penalizes excessive peak temperature, dense line spacing, or long plans.
The forward map $p \mapsto w_{sim}(y; p)$ is evaluated by the calibrated solver. A derivative-free
optimizer (e.g., CMA-ES, Nelder--Mead, or Bayesian optimization) is suitable for the first
implementation.

\subsection{NC Path Generation}
The planned lines are translated into a machine path in a plate-fixed coordinate system. For each
line $i$ at $y_i$, the generator produces start/end points, feed rate, and repetition count
corresponding to the prescribed passes. For a Cartesian gantry, the output can be a G-code program
with linear moves (G01), feedrate changes (F), and torch on/off commands (M-codes). Cooling or
dwell times between passes can be inserted to reflect thermal constraints.

\subsection{Validation Loop}
The resulting NC plan can be re-simulated using the forward model to verify the final curvature
against the target. This closes the loop between planning, execution, and validation, and allows
incremental refinement of the plan parameters.

\end{document}
